\scanpage{195}\mspage{- 60 -}\hfill IX

\uline{i}.\uline{e}. \uline{définit un isomorphisme de chacun de ces schémas abéliens avec le dual de l'autre}.

La démonstration de 5.4 dans le cas général va être donnée au n$^\circ$ 7, comme conséquence d'une construction fort naturelle et importante due à RAYNAUD. Notons seulement ici que pour prouver 5.4, on peut supposer $S$ complet, compte tenu de 3.3.

5.5. Pour exprimer commodément les facteurs extrêmes dans les filtrations (5.1.10) et (5.1.11), dans le cas de la réduction semi-stable, introduisons les schémas en groupes constants tordus sur $S$

\[
\underline{M}'_{\unsure{o|0}}=D(T_o)\ ,\qquad\underline{M}_o=D(T'_o)\ ,\tag*{(5.5.1)}
\]

duals des tores maximaux $T_o$, $T'_o$ des fibres spéciales de $A,A'$ (SGA 3 X 5.7). Comme $S$ est local complet, ces schémas en groupes peuvent être considérés comme les fibres spéciales de schémas constants tordus

\[
\underline{M}'\ ,\ \underline{M}\text{ sur }S\ ;\tag*{(5.5.2)}
\]

ainsi, si le tore $T_o$ est déployé (par exemple si $k$ est séparablement clos) i.e. si $\underline{M}'_o$ est le groupe constant défini par le groupe ordinaire

\[
M'=\mathrm{Hom}_{k\text{-}\mathrm{gr}}(T_o,\mathbb{G}_{mk})\ ,\tag*{(5.5.3)}
\]

alors $\underline{M}'$ est le groupe constant défini par $M'$ :

\[
\underline{M}'=M'_S\ .
\]

On peut d'ailleurs noter que, puisque $T_o$ et $T'_o$ sont isogènes (2.2.6), $T'_o$ est déployé si $T_o$ l'est, de sorte qu'il y a lieu d'introduire également

% archive sheet number: 194
% archive numbering: 378 (pencil and printed label, foot).
